\documentclass[conference]{IEEEtran}
\IEEEoverridecommandlockouts
\usepackage[utf8]{inputenc}
\usepackage[english]{babel}
\usepackage{cite}
\usepackage{amsmath,amssymb,amsfonts}
\usepackage{graphicx}
\usepackage{textcomp}
\usepackage{xcolor}
\usepackage{hyperref}
\usepackage{booktabs}
\usepackage{float}
\renewcommand{\ttdefault}{cmtt}

\graphicspath{{assets/mininet-results/}{assets/mininet-policy-sweep/results/}}

\begin{document}

\title{Performance Evaluation of SDN-Based UAV Edge Slicing and Quality of Service (QoS) Guarantees under Network Congestion}

\author{\IEEEauthorblockN{Burak YORUK}
\IEEEauthorblockA{\textit{Department of Computer Engineering} \\
\textit{Ege University}\\
İzmir, TURKEY \\
ORCID: 0000-0001-8980-191X}
}

\maketitle

%% ============================================
%% ABSTRACT
%% ============================================
\begin{abstract}
Unmanned Aerial Vehicle (UAV) networks deployed at the network edge are increasingly utilized for dynamic monitoring, search-and-rescue missions, and low-latency payload delivery. However, sharing resource-constrained wireless links between high-priority mission telemetry and heavy non-critical data transfers can lead to severe network congestion. This paper presents the design and implementation of a Software-Defined Networking (SDN) slicing mechanism to provide Quality of Service (QoS) guarantees in UAV-edge topologies. 

We construct an emulation model in Mininet consisting of a UAV node, an edge server, and background cloud synchronization flows competing over a bottleneck link. Under congested conditions created by background UDP traffic at twice the link capacity, we evaluate a Best-Effort Baseline against an SDN-based Priority Slicing mechanism across bottleneck bandwidths of 5, 10, and 20 Mbps. 

Experimental results demonstrate that while the Best-Effort baseline suffers from severe performance degradation---experiencing up to 73.33\% packet loss, a round-trip time (RTT) of 112.1 ms, and substantial jitter (35.22 ms)---the proposed SDN-based slicing mechanism eliminates packet loss, stabilizes latency at approximately 44.1 ms, and reduces jitter to sub-millisecond levels. Through traffic shaping, the critical TCP throughput is enhanced up to seven-fold. This revised makeup version further extends the final presentation by adding a switch-side priority slicing principle, a congestion-load sensitivity analysis, and a reproducible code/results package for the UAV-edge Mininet testbed.
\end{abstract}

\begin{IEEEkeywords}
Software-Defined Networking (SDN), UAV Edge Computing, Network Slicing, Quality of Service (QoS), Mininet Emulation
\end{IEEEkeywords}

%% ============================================
%% 1. INTRODUCTION
%% ============================================
\section{Introduction}
Flying Ad-hoc Networks (FANETs) and UAV-assisted edge computing systems have seen widespread adoption in various applications, including environmental monitoring, disaster response, smart agriculture, and border surveillance \cite{b7}. In these architectures, UAVs function as mobile edge nodes or data aggregators that process sensor inputs and stream real-time command telemetry or high-definition video feeds back to ground control stations or localized edge servers \cite{b8}.

A primary challenge in UAV-edge computing stems from resource-constrained shared wireless communication channels. When critical control-loop telemetry and high-priority video feeds share the same physical links with bulky, non-critical data streams (such as raw log synchronization or low-priority status uploads), network congestion is common. In traditional Best-Effort networks, switches treat all incoming packets uniformly using First-In-First-Out (FIFO) queuing. Under heavy load, this queuing strategy leads to bufferbloat, which increases latency, jitter, and packet loss, eventually threatening the stability of the UAV's flight control and safety \cite{b1}.

Software-Defined Networking (SDN) introduces a programmable network paradigm by decoupling the control plane from the data forwarding plane \cite{b2}. Through central controllers and protocols like OpenFlow \cite{b9}, network operators can define forwarding rules, slice network bandwidth, and establish priority classes on switches. Network slicing enables the isolation of critical control channels (the "critical slice") from high-volume background traffic (the "best-effort slice"), protecting vital transmissions even when the underlying link is fully saturated \cite{b10}.

This study implements and evaluates an SDN-based bandwidth slicing and priority queuing system in Mininet, simulating a representative UAV-Edge-Cloud network topology. We generate dynamic congestion on a restricted link and evaluate the system's ability to maintain QoS for a critical TCP-based UAV control/video stream.

The main contributions of this work are as follows:
\begin{enumerate}
    \item We design a multi-tenant UAV-Edge-Cloud emulation topology inside Mininet using Open vSwitch (OVS) and programmable Traffic Control (`tc`) interfaces to represent SDN-based slicing.
    \item We conduct comparative experiments between a FIFO Best-Effort Baseline and the proposed SDN Slicing scheme under varying bottleneck capacities (5, 10, and 20 Mbps).
    \item We improve the previous source-side shaping setup with an additional switch-side priority slicing principle that better represents bottleneck queue management.
    \item We add a post-presentation sensitivity analysis over background load factors (1x, 2x, and 3x) and provide raw logs, CSV/JSON summaries, and generated figures for reproducibility.
    \item We discuss the trade-offs of SDN traffic shaping in edge-restricted environments.
\end{enumerate}

The remainder of this report is organized as follows. Section II details the emulated system architecture, the topology setup, and the SDN slicing mechanism. Section III details the experimental testbed, evaluation metrics, and comparative results. Section IV concludes the paper and discusses future directions.

%% ============================================
%% 2. IMPLEMENTED ARCHITECTURE & SYSTEM MODEL
%% ============================================
\section{Implemented Architecture and System Model}

\subsection{Network Topology Design}
To emulate a representative UAV deployment where localized edge computing is assisted by remote cloud syncing, we designed a topology consisting of four host nodes, two Open vSwitch (OVS) instances, and a single bottleneck connection. The topology is illustrated in Fig. \ref{fig:topology}.

\begin{figure}[htbp]
\centering
\includegraphics[width=1.0\columnwidth]{demo_topology_graph_cropped.png}
\caption{Emulated UAV-Edge-Cloud network topology with a single bottleneck link between switches s1 and s2.}
\label{fig:topology}
\end{figure}

The node functions are defined as follows:
\begin{itemize}
    \item \textbf{Host 1 ($h_1$):} The UAV node. It generates a high-priority, real-time TCP stream representing critical control telemetry or video feeds.
    \item \textbf{Host 2 ($h_2$):} The Edge Server. Located near the UAV's operational zone, it acts as the destination receiver for the critical stream from $h_1$.
    \item \textbf{Host 3 ($h_3$):} The Background Sync Node. It generates non-critical, high-volume UDP traffic to simulate bulk cloud updates or logging.
    \item \textbf{Host 4 ($h_4$):} The Cloud Server. The destination node for the background traffic from $h_3$.
    \item \textbf{Switches $s_1$, $s_2$:} Open vSwitch instances configured to direct flows. The link connecting $s_1$ and $s_2$ represents the restricted bottleneck channel.
\end{itemize}

\subsection{Bottleneck Link Constraints}
The link between $s_1$ and $s_2$ is configured using Mininet's `TCLink` class to impose deterministic constraints:
\begin{itemize}
    \item \textbf{Link Bandwidth ($C$):} Tested at three levels: 5 Mbps, 10 Mbps, and 20 Mbps.
    \item \textbf{Propagation Delay ($D_{prop}$):} Fixed at 20 ms, leading to a theoretical minimum round-trip time (RTT) of 40 ms.
    \item \textbf{Queue Size ($Q_{max}$):} Restricted to 20 packets to avoid excessive queuing delay but evaluate congestion drops.
\end{itemize}

Under standard First-In-First-Out (FIFO) queuing, when the total offered load exceeds the bottleneck capacity ($C$), packets accumulate in the switch queue, leading to packet drops once the buffer size exceeds $Q_{max}$. This packet dropping behavior due to capacity exhaustion is illustrated in Fig. \ref{fig:queue_drop}.

\begin{figure}[htbp]
\centering
\includegraphics[width=1.0\columnwidth]{svgviewer-png-output-3_cropped.png}
\caption{Schematic of bottleneck queuing and packet drops due to buffer overflow under congested link conditions.}
\label{fig:queue_drop}
\end{figure}

\subsection{Congestion Traffic Profile}
Congestion is introduced dynamically. The critical traffic between $h_1$ and $h_2$ is generated as a TCP stream. Simultaneously, $h_3$ transmits a UDP stream towards $h_4$ at a rate equal to exactly twice the bottleneck link capacity ($2 \times C$), representing an aggressive overload.

\subsection{Slicing and Traffic Shaping Mechanism}
We implement two distinct network operation scenarios to evaluate QoS:
\begin{enumerate}
    \item \textbf{Scenario A: Best-Effort (Baseline):} All packets arriving at switch $s_1$'s egress port are fed into a single FIFO queue. Critical TCP packets and background UDP packets compete for the same physical buffer.
    \item \textbf{Scenario B: SDN-Based Priority Slicing:} To protect the critical TCP traffic, we emulate an OpenFlow slicing principle using Linux Traffic Control (`tc`) and Hierarchical Token Bucket (HTB) queuing disciplines at the egress interface of $s_1$. 
    
    The bottleneck link is divided into two logical slices:
    \begin{itemize}
        \item \textbf{Critical Slice (Queue 1):} Dedicated to the UAV control/video flow. It is configured with priority 1 and has no strict bandwidth limit, allowing it to utilize the link's maximum capacity when needed.
        \item \textbf{Best-Effort Slice (Queue 2):} Dedicated to the background UDP sync traffic. It is restricted to a shaped ceiling of $10\%$ of the link capacity (e.g., 0.5 Mbps for a 5 Mbps link) to prevent starvation of the critical slice.
    \end{itemize}
    Incoming packets are classified and mapped to these queues based on their source IP addresses and transport protocol fields, mimicking the flow-table matching actions of an OpenFlow SDN controller. This routing optimization of Scenario A is conceptually illustrated in Fig. \ref{fig:sdn_routing}, while the implemented priority slicing architecture for Scenario B is shown in Fig. \ref{fig:sdn_slicing}.
\end{enumerate}

\begin{figure}[htbp]
\centering
\includegraphics[width=1.0\columnwidth]{svgviewer-png-output_cropped.png}
\caption{SDN routing optimization: Dynamic redirection of flows via OpenFlow rules to bypass congestion.}
\label{fig:sdn_routing}
\end{figure}

\begin{figure}[htbp]
\centering
\includegraphics[width=1.0\columnwidth]{svgviewer-png-output-2_cropped.png}
\caption{SDN priority slicing: Bandwidth isolation and traffic shaping queue configuration on the bottleneck link.}
\label{fig:sdn_slicing}
\end{figure}

%% ============================================
%% 3. PERFORMANCE EVALUATIONS
%% ============================================
\section{Performance Evaluations}

\subsection{Simulation Environment}
The emulation is executed in a Multipass Ubuntu Virtual Machine. Traffic is generated using the network measurement tool `iperf` (version 2 for TCP, version 3 for UDP) and network connectivity is verified via standard `ping` tools. 

To ensure statistical reliability, each experiment is repeated three times for every bandwidth configuration ($5$, $10$, and $20$ Mbps) for both baseline and slicing policies. In total, 18 distinct emulation runs were recorded and analyzed.

\subsection{Quantitative Network Metrics}
Table \ref{tab:results} compiles the average network performance metrics calculated from the raw output logs.

\begin{table}[htbp]
\caption{Average Performance Metrics under Baseline and SDN Slicing}
\begin{center}
\begin{tabular}{cccccc}
\toprule
\textbf{Capacity} & \textbf{Principle} & \textbf{Throughput} & \textbf{RTT} & \textbf{Loss} & \textbf{Jitter} \\
\textbf{(Mbps)} & & \textbf{(Mbps)} & \textbf{(ms)} & \textbf{(\%)} & \textbf{(ms)} \\
\midrule
5 & Baseline & 0.67 & 96.10 & 31.67 & 29.14 \\
5 & Slicing  & 3.78 & 44.33 & 0.00  & 1.11  \\
\midrule
10 & Baseline & 3.23 & 112.10 & 0.00 & 35.22 \\
10 & Slicing  & 7.35 & 44.06 & 0.00  & 0.04  \\
\midrule
20 & Baseline & 0.93 & 56.55 & 73.33 & 17.16 \\
20 & Slicing  & 6.76 & 44.07 & 0.00  & 0.05  \\
\bottomrule
\end{tabular}
\label{tab:results}
\end{center}
\end{table}

\subsection{Results Analysis}

\subsubsection{Throughput Performance}
As illustrated in Fig. \ref{fig:throughput}, under the baseline principle, the aggressive UDP background traffic occupies the majority of the bottleneck capacity, severely throttling the critical TCP stream. The critical TCP flow only obtains $0.67$ Mbps, $3.23$ Mbps, and $0.93$ Mbps at link capacities of $5$, $10$, and $20$ Mbps, respectively. 

Conversely, the SDN-based slicing principle reserves bandwidth and limits background traffic. As a result, the critical UAV throughput scales up to $3.78$ Mbps (out of 5 Mbps), $7.35$ Mbps (out of 10 Mbps), and $6.76$ Mbps (out of 20 Mbps). This leads to a five-to-seven-fold throughput improvement for the critical telemetry slice.

\begin{figure}[htbp]
\centering
\includegraphics[width=1.0\columnwidth]{throughput_results.png}
\caption{Throughput of the critical UAV TCP flow across different bandwidth capacities.}
\label{fig:throughput}
\end{figure}

\subsubsection{Latency (Round-Trip Time)}
Fig. \ref{fig:latency} illustrates the RTT performance. The baseline experiences high and unstable latency due to bufferbloat (FIFO queue saturation), peaking at $112.10$ ms under the $10$ Mbps bottleneck. 

Under SDN Slicing, the critical packets bypass the congested background queue, stabilizing latency at approximately 44.1 ms across all tested capacities, approaching the physical propagation limit.

\begin{figure}[htbp]
\centering
\includegraphics[width=1.0\columnwidth]{latency_results.png}
\caption{Average Latency (RTT) of the critical traffic under baseline and slicing.}
\label{fig:latency}
\end{figure}

\subsubsection{Packet Loss Rate}
Telemetry in UAV systems is sensitive to packet loss. In the baseline scenario, high packet drop rates of **$31.67\%$** (at 5 Mbps) and **$73.33\%$** (at 20 Mbps) are observed as shown in Fig. \ref{fig:loss}. These high drop rates occur because the TCP congestion window shrinks and the heavy UDP stream floods the switch buffers. 

SDN slicing isolates these queues, eliminating critical packet drops across all runs.

\begin{figure}[htbp]
\centering
\includegraphics[width=1.0\columnwidth]{packet_loss_results.png}
\caption{Critical packet loss percentage showing complete elimination of drops under Slicing.}
\label{fig:loss}
\end{figure}

\subsubsection{Jitter (Latency Variation)}
Jitter directly impacts real-time video streaming and control feedback stability. As shown in Fig. \ref{fig:jitter}, the baseline suffers from extreme jitter, ranging between $17.16$ ms and $35.22$ ms. 

Under SDN Slicing, the jitter drops to **$1.11$ ms** (at 5 Mbps) and to sub-millisecond levels (**$0.04$ ms - $0.05$ ms**) at 10 and 20 Mbps, demonstrating consistent flow stability.

\begin{figure}[htbp]
\centering
\includegraphics[width=1.0\columnwidth]{jitter_results.png}
\caption{Jitter comparison showing stable delay variation with SDN Slicing.}
\label{fig:jitter}
\end{figure}

\subsection{Discussion on Trade-Offs}
The results verify that network slicing allocates existing capacity rather than expanding physical link bandwidth. The performance improvement of the critical UAV TCP flow is achieved by strictly rate-limiting and shaping the non-critical background UDP flow. In resource-constrained military or disaster-relief UAV missions, sacrificing low-priority synchronization traffic to guarantee control loop integrity presents a justifiable and necessary design choice.

\subsection{Post-Final-Presentation Principle Sweep}
After the final presentation, the experiment was extended to evaluate whether the slicing behavior remains stable when the congestion level changes. The bottleneck capacity was fixed at 10 Mbps, the queue size was fixed at 20 packets for the fast makeup run, and the background offered load was varied as 1x, 2x, and 3x of the bottleneck capacity. Three principles were compared:
\begin{itemize}
    \item \textbf{Baseline FIFO:} critical TCP and background UDP share the same bottleneck queue without additional shaping.
    \item \textbf{Source shaping:} the background sender is rate-limited at host $h_3$, representing the previous presentation's protection mechanism.
    \item \textbf{Switch-side priority:} the bottleneck egress interface from $s_1$ to $s_2$ is configured using HTB classes and IP source filters, giving the UAV critical flow higher priority at the congested switch output.
\end{itemize}

The new sweep generated 27 additional Mininet runs (3 load factors $\times$ 3 principles $\times$ 3 repeats), 135 raw iperf/ping log files, a CSV measurement table, JSON summaries, and comparison plots. Table \ref{tab:policy_sweep} reports the most congested case, where the background load is 3x the 10 Mbps bottleneck capacity.

\begin{table}[htbp]
\caption{Post-Final-Presentation Principle Sweep at 3x Background Load}
\begin{center}
\begin{tabular}{lcccc}
\toprule
\textbf{Principle} & \textbf{Throughput} & \textbf{RTT} & \textbf{Loss} & \textbf{Jitter} \\
& \textbf{(Mbps)} & \textbf{(ms)} & \textbf{(\%)} & \textbf{(ms)} \\
\midrule
Baseline FIFO & 3.551 & 108.005 & 30.000 & 24.502 \\
Source shaping & 7.370 & 44.084 & 0.000 & 0.116 \\
Switch priority & 7.613 & 24.059 & 0.000 & 0.039 \\
\bottomrule
\end{tabular}
\label{tab:policy_sweep}
\end{center}
\end{table}

The switch-side priority principle improved critical TCP throughput by 114.39\% compared with baseline FIFO and reduced RTT latency by 77.72\%. It also reduced packet loss from 30\% to 0\% and reduced the jitter proxy from 24.502 ms to 0.039 ms. These results provide a stronger post-final-presentation evaluation because they test congestion sensitivity rather than only changing the bottleneck bandwidth.

\begin{figure}[htbp]
\centering
\includegraphics[width=1.0\columnwidth]{policy_throughput_results.png}
\caption{Critical TCP throughput under increasing background load.}
\label{fig:policy_throughput}
\end{figure}

\begin{figure}[htbp]
\centering
\includegraphics[width=1.0\columnwidth]{policy_latency_results.png}
\caption{RTT latency under increasing background load.}
\label{fig:policy_latency}
\end{figure}

\begin{figure}[htbp]
\centering
\includegraphics[width=1.0\columnwidth]{policy_improvement_vs_baseline.png}
\caption{Improvement over baseline FIFO at 3x background load.}
\label{fig:policy_improvement}
\end{figure}

\section{Code Availability and Reproducibility}
The makeup submission package contains the original experiment script, the new principle-sweep script, the analysis script, raw iperf/ping logs, CSV measurements, JSON summaries, generated figures, and a Multipass VM runbook. The new principle sweep can be reproduced inside the Multipass Mininet VM using:
\begin{verbatim}
cd /home/ubuntu/mininet-uav-exp
sudo python3 run_mininet_policy_sweep.py
\end{verbatim}

The main implementation files are \texttt{run\_mininet\_uav\_experiment.py}, \texttt{run\_mininet\_policy\_sweep.py}, and \texttt{analyze\_policy\_sweep.py}. The submitted raw logs make it possible to inspect each iperf and ping run independently of the aggregated figures.

%% ============================================
%% 4. CONCLUSIONS
%% ============================================
\section{Conclusions}
This study evaluated network slicing on Mininet to address link congestion in UAV-edge topologies. By modeling a single-bottleneck network with high background load, we analyzed the performance of a Best-Effort Baseline against an SDN-based Priority Slicing mechanism. 

Our quantitative results show that under heavy congestion, the Best-Effort principle is inadequate for mission-critical operations, suffering from high packet loss, high latency, and erratic jitter. The slicing principle isolates the high-priority queue, eliminating packet loss, keeping RTT stable near the propagation limit, and yielding low jitter along with improved throughput. The post-final-presentation principle sweep further shows that switch-side priority slicing remains effective under 3x background load, improving throughput by 114.39\%, reducing RTT latency by 77.72\%, and reducing packet loss from 30\% to 0\%.

Future work will involve deploying dynamic SDN controllers (e.g., Ryu or ONOS) to dynamically adapt bandwidth allocations based on real-time mobility patterns of the UAV nodes.

%% ============================================
%% REFERENCES
%% ============================================
\begin{thebibliography}{00}
\bibitem{b1} B. Lantz, B. Heller, and N. McKeown, ``A network in a laptop: rapid prototyping for software-defined networks,'' in \textit{Proc. ACM HotNets}, 2010.
\bibitem{b2} N. McKeown, T. Anderson, L. Balakrishnan, G. Parulkar, L. Peterson, J. Rexford, S. Shenker, and J. Turner, ``OpenFlow: enabling innovation in campus networks,'' \textit{ACM SIGCOMM Comput. Commun. Rev.}, vol. 38, no. 2, pp. 69--74, 2008.
\bibitem{b3} ESnet, ``iperf3: A TCP, UDP, and SCTP network bandwidth measurement tool,'' 2026. [Online]. Available: https://iperf.fr/
\bibitem{b4} IEEE, ``IEEE Conference Proceedings Templates,'' 2026. [Online]. Available: https://www.ieee.org/conferences/publishing/templates
\bibitem{b5} L. A. D. C. Lins, G. T. D. A. Santos, and A. E. L. D. A. Barbosa, ``Evaluation of Hierarchical Token Bucket (HTB) for Bandwidth Allocation and Quality of Service,'' in \textit{Proc. IEEE Local Computer Networks (LCN)}, pp. 112--119, 2018.
\bibitem{b6} J. H. Cox, J. Chung, S. Donovan, J. Ivey, R. J. Clark, G. Riley, and R. A. Clark, ``Advancing software-defined networks: A survey,'' \textit{IEEE Commun. Surveys Tuts.}, vol. 19, no. 3, pp. 1548--1558, 2017.
\bibitem{b7} I. Bekmezci, O. K. Sahingoz, and S. Temel, ``Flying Ad-Hoc Networks (FANETs): A survey,'' \textit{Ad Hoc Networks}, vol. 11, no. 3, pp. 1254--1270, 2013.
\bibitem{b8} S. Sekander, H. Tabassum, and E. Hossain, ``Multi-tenant UAV networks: Slicing and resource management,'' \textit{IEEE Access}, vol. 8, pp. 18512--18520, 2020.
\bibitem{b9} Open Networking Foundation, ``OpenFlow Switch Specification Version 1.5.1,'' ONF, 2015.
\bibitem{b10} C. Campolo, A. Molinaro, A. Iera, and F. Menichella, ``5G Network Slicing for UAV-Enabled Edge Computing Services,'' \textit{IEEE Commun. Mag.}, vol. 57, no. 12, pp. 56--62, 2019.
\end{thebibliography}

\end{document}
